%! TEX program = xelatex
\documentclass[a4paper,11pt]{report}

\usepackage{./../packages/mainstyle}
\usepackage{./../packages/colors}
\usepackage{./../packages/frameboxes}
\usepackage{./../packages/title}
\usepackage{./../packages/packs}
\usepackage{./../packages/macros}

\usepackage{tabularx} 
\usepackage{array} 
\usepackage[table]{xcolor}

\usepackage[italian]{babel}

\usepackage[T1]{fontenc}
\usepackage{FiraMono} 

\usepackage{listings}
\usepackage{xcolor}

\tcbuselibrary{skins, listings}

\definecolor{hsbg}{HTML}{F8F9FA}       
\definecolor{hspurple}{HTML}{5E5086}   
\definecolor{hsgreen}{HTML}{118C4F}   
\definecolor{hsred}{HTML}{C0392B}      
\definecolor{hsnumber}{HTML}{ADB5BD}  
\definecolor{boxtitlebg}{HTML}{2C3E50} 

\lstdefinestyle{haskellstyle}{
    language=Haskell,
    commentstyle=\color{hsgreen}\itshape,
    keywordstyle=\color{hspurple}\bfseries,
    numberstyle=\tiny\color{hsnumber},
    stringstyle=\color{hsred},
    basicstyle=\ttfamily\footnotesize, 
    breaklines=true,
    showstringspaces=false,
    tabsize=2,                        
    numbers=none                     
}

% for few lines of code
\newtcblisting{snippet}{
    enhanced,
    colback=hsbg,
    colframe=hspurple,                 
    boxrule=0pt,                      
    leftrule=3pt,                    
    arc=0pt,                        
    outer arc=0pt,
    listing only,
    listing options={style=haskellstyle},
    left=8pt,                        
    top=2pt,
    bottom=2pt,
    boxsep=0pt,
    before skip=6pt,                
    after skip=6pt
}

% for large amounts of code
\newtcblisting{haskellbox}[1]{
    enhanced,
    skin=bicolor,
    colback=hsbg,
    colbacktitle=boxtitlebg,
    coltitle=white,
    colframe=boxtitlebg,
    boxrule=0.5pt,
    arc=4pt,
    drop shadow=black!10!white,
    fonttitle=\bfseries\sffamily,
    title=#1,
    listing only,
    listing options={style=haskellstyle, numbers=left, numbersep=10pt}, % Add numbers here!
    left=15pt,
    top=4pt,
    bottom=4pt,
    breakable
}

% for inline code

\definecolor{inlinebg}{HTML}{F3F4F6}    
\definecolor{inlinetext}{HTML}{24292E} 

\newtcbox{\code}{
    on line,                           
    fontupper=\ttfamily\color{inlinetext},
    colback=inlinebg,                 
    colframe=inlinebg,               
    boxrule=0pt,                    
    arc=2pt,                       
    boxsep=0pt,                   
    left=3pt,                    
    right=3pt,
    top=1pt,                    
    bottom=1pt
}

\usepackage[autostyle=false]{csquotes}
\MakeOuterQuote{"}

\usepackage{float}
\usepackage{enumitem}
\usepackage{bussproofs}
\usepackage{mdframed}

\usepackage{titlesec}

\setmainfont{XCharter}

\setcoursename{Tecniche di Programmazione\\ Funzionale e Imperativa}
\setcoursebook{slide del professor Salvo}
\setauthorname{aglaia norza}
\setauthoremail{thisisaglaia@gmail.com}
\setauthorgithub{AglaiaNorza}
\setcoverphoto{covermeme.png}

\setcoverwidth{0.6\textwidth}
\setcoverheight{0.6\textwidth}
\setbeforespace{2.5em}
\setafterspace{2.5em}

\begin{document}

\makecustomcover

\tableofcontents
\newpage

\chapter{Le basi}
\section{Concetti fondamentali}
\begin{itemize}
    \item \textbf{Valutazione lazy:} Haskell valuta le espressioni solo quando è strettamente necessario (permette di gestire agevolmente strutture dati infinite)
    \item \textbf{Trasparenza referenziale:} non essendoci assegnazioni o memoria mutabile (assenza di side-effects), una variabile identifica sempre lo stesso oggetto (permette di manipolare algebricamente i programmi)
    \item \textbf{Polimorfismo e type inference:} molte funzioni operano su tipi generici (es. \code{a}, \code{b}). Il compilatore inferisce automaticamente il \textit{tipo principale} (il più generale possibile); tutti gli altri tipi corretti sono sue istanze.
\item \textbf{Non-terminazione, Bottom e \code{undefined}}: in Haskell, ogni tipo possiede implicitamente un valore speciale (chiamato \textit{bottom}, indicato matematicamente con $\bot$) che rappresenta un'eccezione fatale o una computazione che entra in un loop infinito
    \begin{itemize}
        \item \textbf{la costante \code{undefined}:} è la materializzazione a livello di codice del valore $\bot$ - poiché una computazione che fallisce o non termina può avvenire in qualsiasi contesto (mentre si aspetta un numero, un booleano, ecc.), \code{undefined} possiede il tipo più generale e polimorfo possibile: \code{a} (ovvero $\forall \alpha . \alpha$)
        \item \textbf{Il limite del lambda calcolo (\textit{Occurs check}):} nel lambda calcolo puro, il loop infinito per eccellenza è il combinatore $\Omega = (\lambda x.xx)(\lambda x.xx)$. In Haskell, l'auto-applicazione $\lambda x.xx$ \textbf{non è tipabile}. Perché $x$ possa prendere se stesso come argomento, il suo tipo $\alpha$ dovrebbe essere uguale a una funzione che prende $\alpha$ come parametro e restituisce $\beta$ ($\alpha = \alpha \to \beta$). Questo creerebbe un tipo infinito (durante l'inferenza dei tipi, il compilatore esegue un controllo chiamato \textit{Occurs check} che impedisce alla variabile di tipo $\alpha$ di apparire all'interno della sua stessa definizione, rifiutando di conseguenza l'espressione)
        \item \textbf{Ricorsione esplicita:} poiché il type system ci vieta di creare loop tramite l'auto-applicazione anonima di $\lambda x.xx$, in Haskell la non-terminazione si esprime unicamente attraverso la ricorsione esplicita. Definendo ad esempio \code{omega' = omega'}, creiamo intenzionalmente una computazione divergente che il compilatore accetta, assegnandole il tipo più generale \code{a}.
    \end{itemize}
    \item \textbf{Ideologia funzionale:} l'ideologia funzionale favorisce l'assemblaggio di funzioni semplici tramite composizione (\code{.}). Si pensa in termini di trasformazioni globali di strutture dati, per poi raffinare in un secondo momento i colli di bottiglia prestazionali.
\end{itemize}

\section{Sintassi, Funzioni e Pattern Matching}

\subsection{Definizione e Applicazione}
\begin{itemize}
    \item l'applicazione di funzione: \textbf{associa a sinistra} e si scrive senza virgole (\code{f x y}).
    \item \textbf{aperatori infissi vs prefissi:} le funzioni binarie si usano in modo infisso tra backtick \\ ( \code{10 `div` 2}), gli operatori infissi in modo prefisso tra parentesi ( \code{(+) 3 4}).
\end{itemize}

\subsection{Lambda Espressioni e Combinatori}
Le funzioni anonime si scrivono con backslash (\code{\textbackslash} simula $\lambda$).
\begin{haskellbox}{}
-- identita' (combinatore I)
i :: t -> t
i x = x           -- definizione standard
i' = \x -> x      -- lambda espressione

-- combinatore K (proiettore del primo argomento) (primo assioma di Hilbert)
k :: t1 -> t2 -> t1
k' = \x y -> x

-- combinatore S (sostituzione) (secondo assioma di Hilbert)
s :: (t2 -> t1 -> t) -> (t2 -> t1) -> t2 -> t
s x y z = x z (y z)
\end{haskellbox}

\subsection{Pattern matching}
Il pattern matching decompone le strutture dati in base alla loro forma sintattica. L'ordine delle clausole è strettamente sequenziale (viene scelta la prima dall'alto che fa matching).
\begin{itemize}
    \item \textbf{mancanza di unificazione:} non si possono usare variabili identiche nel pattern per verificarne l'uguaglianza (es. \code{f x x = ...} è invalido)
    \item \textbf{underscore (\code{\_}):} variabile anonima ("don't care") per valori non rilevanti
    \item \textbf{as-pattern (\code{@}):} permette di dare un nome all'intera struttura mentre la si decompone (es. \code{xs@(x:txs)})
\item \textbf{guardie (\code{|}):} permettono di definire il comportamento di una funzione valutando sequenzialmente una serie di \textit{predicati booleani} (condizioni logiche); a differenza del pattern matching che controlla la \textit{forma} sintattica dei dati, le guardie valutano \textit{valori e relazioni} (es. \code{x > 0}, \code{x == y})
    \begin{itemize}
    \item l'ultima guardia è tipicamente \code{otherwise} (che nel Prelude di Haskell è semplicemente un alias per \code{True}), che funge da caso di default ("catch-all") per assicurarsi che la funzione restituisca sempre un valore.
    \end{itemize}
\begin{haskellbox}{}
confronta x y
  | x == y    = "i numeri sono uguali" -- supera la mancanza di unificazione
  | x > y     = "il primo e' maggiore"
  | otherwise = "il secondo e' maggiore"
\end{haskellbox}{}
\item \textbf{equazioni multiple (definizione per casi):} il modo più comune di fare pattern matching è definire la funzione scrivendo più "intestazioni" separate; il compilatore le prova dall'alto verso il basso: se gli argomenti combaciano con la forma (pattern) della riga, esegue quell'equazione, altrimenti controlla la successiva.
\begin{haskellbox}{}
    myAnd True True = True
    myAnd   _   _   = False
\end{haskellbox}{}
\begin{itemize}
    \item è meno comune, ma è possibile fare pattern-matching anche nelle \textit{lambda-espressioni}:
        \begin{haskellbox}{}
            myFst' = \(x, y) -> x
        \end{haskellbox}
\end{itemize}
\end{itemize}

\subsection{Scope locale: \code{let} e \code{where}}
\begin{itemize}
    \item \textbf{\code{let ... in}:} costrutto standalone per definire variabili/funzioni locali (utile per strutturare/destrutturare tuple nelle ricorsioni)
    \item \textbf{\code{where}:} clausola che qualifica la parte destra di una definizione appena conclusa
\end{itemize}
\begin{haskellbox}{}
fibAux n = let (f, fPrec) = fibAux (n-1) in (f + fPrec, f)
-- equivalente a:
fibAux' n = (f + fPrec, f) where (f, fPrec) = fibAux' (n-1)
\end{haskellbox}

\begin{gframe}{differenza tra \code{let} e \code{where}}
   una clausola \code{where} è visibile da tutte le guardie di una specifica equazione. Un \code{let} definito dopo il simbolo \code{=} di una guardia è visibile solo all'interno di quella singola guardia. Per condividere un calcolo tra più guardie, il \code{where} è obbligatorio.
\end{gframe}

\begin{haskellbox}{}
-- condivisione tra guardie
analizzaRaggio r
  | area > 100 = "cerchio grande"
  | area < 10  = "cerchio piccolo"
  | otherwise  = "cerchio medio"
  where area = pi * r^2  -- 'area' calcolata una volta e visibile a tutti i pipe

-- inline: (serve il let perche' non c'e' una dichiarazione di funzione)
volumeCilindro r h = (let area = pi * r^2 in area) * h
\end{haskellbox}

\newpage
\section{Sistema dei Tipi e Classi}

Haskell è \textbf{statically type-safe}: se un programma è tipato correttamente, nessun errore di tipo si verifica durante l’esecuzione.

\subsection{Type Signatures}
In Haskell, grazie al meccanismo di type inference, il compilatore è in grado di dedurre automaticamente il tipo più generale per quasi ogni espressione. Tuttavia, è considerata un'ottima pratica (e in alcuni casi più complessi è strettamente necessario) dichiarare esplicitamente il tipo di una funzione subito sopra la sua definizione. 

Questo si fa usando l'operatore \code{::} (che si legge "ha tipo"). Dichiarare esplicitamente i tipi è utile ai fini della leggibilità del codice e permette di "forzare" il tipo desiderato, aiutando il compilatore a individuare eventuali errori logici in modo molto più preciso.

\begin{haskellbox}{}
-- dichiarazione esplicita del tipo: prende due Int e restituisce un Int
somma :: Int -> Int -> Int
somma x y = x + y
\end{haskellbox}

\subsection{Type Classes}
Le classi limitano il polimorfismo imponendo vincoli (o "interfacce") sui tipi. Più che semplici insiemi, rappresentano classificazioni di tipo algebrico: raggruppano i tipi su cui sono definite certe operazioni. Le classi possono formare gerarchie, dove una sottoclasse "estende" la superclasse ereditandone le operazioni e aggiungendone di nuove.

\begin{itemize}
    \item \textbf{Eq (Equality Types):} comprende i tipi che ammettono l'uguaglianza (\code{==}) e la sua negazione (\code{/=}). Tutti i tipi base (Bool, Int, Float, ecc.) sono istanze di Eq, così come le liste e tuple costruite su tipi Eq
        \subitem \textit{definizione interna di Eq}
        \begin{haskellbox}{}
class Eq a where
  (==), (/=) :: a -> a -> Bool
  -- le classi possono contenere codice (metodi di default)
  x /= y = not (x == y)
\end{haskellbox}

    \item \textbf{Ord (Ordered Types):} sottoclasse di \code{Eq} - tipi ordinabili che supportano operazioni come \code{<=}, \code{<}, \code{>=}, \code{>}, \code{min} e \code{max}.
\begin{itemize}
    \item per definire un'istanza di Ord, è sufficiente fornire il codice per l'operatore \code{<}, il resto può essere interdefinito.
\end{itemize}

    \item \textbf{Num:} la classe più generica per i tipi numerici. È sottoclasse di \code{Eq} e \code{Show}. Implementa operazioni aritmetiche base (\code{+}, \code{-}, \code{*}), oltre a valore assoluto (\code{abs}) e segno (\code{signum}).
        \subitem i numeri interi costanti sono polimorfi rispetto alla classe \code{Num}:
        \begin{haskellbox}{}
> :t 3
3 :: Num a => a  -- 3 puo' essere un Int, Integer, Float, ma non un Char
\end{haskellbox}

    \item \textbf{Altre classi utili:}
        \begin{itemize}
            \item \textbf{Show:} tipi che possono essere stampati (convertiti in \code{String}) tramite il metodo \code{show}. Se provate a stampare a video il risultato di una funzione che non appartiene a \code{Show} (es. una lambda anonima come \code{\textbackslash x -> x}), l'interprete darà errore.
            \item \textbf{Read:} permette di trasformare una \code{String} in un valore tramite il metodo \code{read}
            \item \textbf{Integral:} sottoclasse di \code{Num} per i numeri interi (es. Int, Integer); aggiunge le operazioni di divisione intera \code{div} e modulo \code{mod}
            \item \textbf{Fractional:} sottoclasse di \code{Num} per i numeri frazionari (es. Float, Rational); aggiunge l'operatore di divisione reale \code{(/)} e \code{recip}
        \end{itemize}
\end{itemize}

\subsubsection{Vincoli di classe}
I vincoli di classe si indicano prima del tipo vero e proprio, separati dalla doppia freccia \code{=>}. Vincoli multipli si racchiudono tra parentesi e indicano che la funzione opera su un tipo che deve soddisfare tutte le classi elencate.

\begin{haskellbox}{}
-- esempio: mcd richiede che il tipo 'a' supporti sia l'ordine (<, ==) 
-- sia le operazioni aritmetiche (-)
mcd :: (Ord a, Num a) => a -> a -> a
mcd x y
  | x == y    = x
  | x < y     = mcd (y - x) x
  | otherwise = mcd (x - y) y
\end{haskellbox}

\subsection{Definizione di Tipi: \code{type} e \code{data}}
\begin{itemize}
    \item \textbf{\code{type} (sinonimi):} danno nuovi nomi a tipi esistenti (es. \code{type Casella = (Int, Int)})   
    \item \textbf{\code{data} (tipi algebrici):} costruiscono nuovi tipi tramite costruttori finiti, parametrici o ricorsivi. Possono essere dichiarati istanze di classi.
\end{itemize}

\begin{haskellbox}{}
data SetteNani = Eolo | Pisolo | Brontolo          -- finiti
data Either a b = Left a | Right b                 -- unioni disgiunte
data Maybe a = Just a | Nothing                    -- gestione fallimenti
data List a = Nil | Cons a (List a)                -- ricorsivi/induttivi
data BTree a = Foglia a | Radice (BTree a) (BTree a)
\end{haskellbox}
\subsubsection{Il costruttore di tipo \code{Maybe}}
Il tipo \code{Maybe} è il tipo di computazioni che possono fallire. In Haskell, è il tipo con cui si rappresentano comunemente le eccezioni in maniera sicura e controllata dal sistema di tipi (corrisponde a \code{optional})

Un valore di tipo \code{Maybe} può assumere due forme:
\begin{itemize}
    \item \code{Just v}: contiene un valore \code{v} (la computazione ha avuto successo e restituisce qualcosa)
    \item \code{Nothing}: non contiene nulla (rappresenta una computazione indefinita o fallita)
\end{itemize}

Si può utilizzare per le \textit{funzioni parziali} (che non sono definite per tutti i possibili valori di input (eg divisione per zero)). Tramite il pattern matching, possiamo "spacchettare" il dato per estrarne il contenuto.

\begin{haskellbox}{}
-- definizione del tipo predefinito
data Maybe a = Just a | Nothing

-- 1. costruzione di un Maybe: divisione sicura
-- Se il divisore e' 0 restituisce Nothing, altrimenti restituisce Just il risultato
safediv :: Integral a => a -> a -> Maybe a
safediv n 0 = Nothing
safediv n m = Just (n `div` m)

-- 2. destrutturare un Maybe: funzione "inversa" per estrarre il valore
extract :: Maybe a -> a
extract (Just n) = n
extract Nothing  = undefined 

-- extract (safediv 30 0) -> *** Exception: Prelude.undefined 
\end{haskellbox}

\subsection{Istanze e Derivazione (\code{deriving})}
La keyword \code{deriving} è un meccanismo automatico che permette al compilatore di generare le istanze per le classi predefinite senza dover scrivere codice manuale. 

\begin{itemize}
    \item \textbf{Eq:} genera un'uguaglianza basata sulla sintassi
    \item \textbf{Ord:} deriva l'ordine lessicografico in base alla sequenza dei costruttori
    \item \textbf{Show/Read:} la conversione usa i nomi dei costruttori come stringhe
    \item \textbf{Enum:} permette di generare liste per enumerazione (es. \code{[Lun..Ven]})
\end{itemize}

\begin{haskellbox}{}
-- istanze automatiche tramite deriving 
data Giorni = Lun | Mar | Mer | Gio | Ven | Sab | Dom
    deriving (Eq, Ord, Show, Enum, Read)

-- esempio di istanza manuale 
instance Eq Bool where
    False == False = True
    True == True   = True
    _ == _         = False
\end{haskellbox}

\subsection{Tipabilità e Let-Polymorphism}
\begin{itemize}
    \item Non tutto è tipabile in Haskell. Il classico esempio è l'auto-applicazione pura:
    \begin{itemize}
        \item \code{omega = \textbackslash x -> x x} non è tipabile perché richiederebbe un tipo infinito ($t = t \to t_1$), portando a una teoria dei tipi non decidibile (e quindi a tipi non inferibili dal compilatore)
    \end{itemize}
    
    \item Dal punto di vista del calcolo, \code{let x = N in M} è equivalente all'applicazione di una lambda: \code{(\textbackslash x -> M) N}. Tuttavia, il compilatore Haskell tratta i due casi in modo diverso durante l'inferenza dei tipi.

    \item \textbf{Let-Polymorphism:}

    Il costrutto \code{let} permette di assegnare a una variabile un tipo polimorfo che viene "generalizzato" prima di essere usato nel corpo (\code{in}).
    \begin{itemize}
        \item \textbf{tipabile:} \code{let x = \textbackslash y -> y in x x} è tipabile.

            Haskell vede che \code{x} è l'identità ($a \to a$) e, nel corpo \code{x x}, può istanziare il primo \code{x} come $(a \to a) \to (a \to a)$ e il secondo come $(a \to a)$.
        \item \textbf{non tipabile:} \code{(\textbackslash x -> x x) (\textbackslash y -> y)} \textbf{non} è tipabile. 

            Qui il compilatore deve tipare la funzione \code{\textbackslash x -> x x} prima di sapere che \code{x} sarà l'identità. Come visto sopra, \code{\textbackslash x -> x x} è intrinsecamente non tipabile.
    \end{itemize}

\end{itemize}

\begin{lemmaframe}{Terminazione}{}
    Tutti i lambda-termini tipabili nel sistema dei tipi di Haskell (che non usino ricorsione esplicita tramite equazioni) \textbf{terminano}.
\end{lemmaframe}
 

\subsection{Semantica Denotazionale: ricorsione e punti fissi}

Haskell è un linguaggio dichiarativo: una definizione non è un comando, ma un'equazione. 
Quando definiamo una funzione ricorsiva $f$, stiamo scrivendo un'equazione del tipo $f = \mathcal{E}(f)$, dove $\mathcal{E}$ è un'espressione che contiene $f$ stessa.

\begin{itemize}
    \item \textbf{Bottom ($\bot$)} rappresenta il "non-valore". Indica un calcolo che non termina (loop infinito) o che fallisce (errore). In Haskell, ogni tipo $t$ contiene implicitamente $\bot$.
    \item Se la definizione è un'equazione, la "semantica" (ovvero la funzione reale prodotta dal compilatore) deve essere una soluzione dell'equazione. Matematicamente, una soluzione di $f = T(f)$ è un \textbf{punto fisso} del funzionale $T$.
\end{itemize}

\begin{gframe}{Approfondimento: ricorsione e non-terminazione: \code{omega'}}
Prendiamo la funzione
\begin{haskellbox}{}
omega' x = omega' x
-- > :t omega'
-- omega' :: t -> t1
\end{haskellbox}
Questa funzione ha il tipo più generico possibile (\code{t -> t1}) perché, non terminando mai, non produce mai un valore reale.

Rimuovendo $x$, otteniamo $omega' = I(omega')$, dove $I$ è l'identità.

Qualsiasi funzione $f$ soddisfa $f = I(f)$. Quale scegliere?

\begin{itemize}
    \item Haskell sceglie la funzione che "fa il minimo indispensabile" per soddisfare l'equazione. Si usa l'ordinamento $\sqsubseteq$ dove $f \sqsubseteq g$ significa che $g$ è "più definita" di $f$.
    \item \textbf{Costruzione di Kleene (least-fixpoint):} 

        il computer "costruisce" la funzione partendo dal nulla ($\bot$):
    \[ f_0 = \bot \quad (\text{non so nulla}) \]
    \[ f_1 = T(f_0) \quad (\text{applico la definizione una volta}) \]
    \[ f_{\infty} = \text{limite della catena} = \textit{lfp}(T) \]
    Nel caso di \code{omega'}, la catena è $\{\bot, \bot, \dots\}$. Il limite è $\bot$. La semantica di \code{omega'} è dunque la funzione che restituisce sempre "non-valore".
\end{itemize}

\end{gframe}

Il punto fisso è quindi la "semantica" perché è l'unico modo per dare un valore univoco e calcolabile a una definizione circolare. Il \textit{minimo} punto fisso garantisce che il significato sia quello "prodotto" effettivamente dal calcolo ricorsivo.

\begin{gframe}{\code{myundefined}}
Mentre \code{omega'} è una \textit{funzione} che non termina, possiamo definire un \textit{valore} che non termina:
\begin{haskellbox}{}
myUndefined = myUndefined
-- > :t myUndefined
-- myUndefined :: a
\end{haskellbox}

\begin{itemize}
    \item  Poiché \code{myUndefined} ha tipo \code{a}, esso \textbf{appartiene a tutti i tipi}. È l'implementazione "home-made" della costante \code{undefined} di Haskell. Rappresenta il valore $\bot$ per ogni possibile tipo del linguaggio.
    \item Matematicamente, \code{myUndefined} è il \textbf{punto fisso dell'identità}: 
\begin{enumerate}
\item la definizione \code{x = x} corrisponde all'equazione funzionale $x = I(x)$
\item la soluzione di questa equazione è, per definizione, il punto fisso dell'identità $I$
\item poiché Haskell usa la semantica del minimo punto fisso, e l'identità applicata al valore nullo ($\bot$) restituisce ancora $\bot$, il risultato è la non-terminazione
\item questo spiega perché \code{myUndefined} ha tipo \code{a}: non essendo "sceso" verso nessun valore specifico (come un intero o una stringa), rimane una pura astrazione di errore/loop che può "fingere" di essere qualsiasi tipo.
\end{enumerate}
\end{itemize}



Haskell permette quindi di manipolare la non-terminazione come un valore reale. Ogni tipo in Haskell è in realtà "puntato" (\textit{pointed}), ovvero contiene sempre almeno un valore: $\bot$.
\end{gframe}




\section{Strutture Dati Predefinite}

\subsection{Tipi di Base e Tuple}
\begin{itemize}
    \item \textbf{Booleani:} \code{Bool} (\code{True}, \code{False}). Operatori: \code{not}, \code{\&\&}, \code{||}.
    \item \textbf{Interi:} \code{Int} (dimensione fissa), \code{Integer} (precisione illimitata).
    \item \textbf{Tuple:} n-uple fisse e disomogenee (es. \code{(String, Bool)})
        \subitem offrono le funzioni \textit{su coppie}: \code{fst}, \code{snd} (prendono primo e secondo argomento)
\end{itemize}

\subsection{Liste}
Sequenze di elementi omogenei. Definite dal costruttore vuoto \code{[]} e dall'operatore infisso \textit{cons} \code{(:)}. Le stringhe sono \code{[Char]}. 

Zucchero sintattico: \code{[1,2,3]} equivale a \code{1:2:3:[]}.

\begin{haskellbox}{}
-- enumerazioni 
[m..n]           -- range finito (lista con elementi da m a n)
[m, n..p]        -- range con step (es. [1, 3..11] dispari)
[m..]            -- lista infinita (lazy)

\end{haskellbox}
Haskell supporta le \textbf{list comprehension} (definizioni set-like del contenuto di una lista)

\begin{haskellbox}{}
-- list comprehension (spesso tradotte internamente in map/filter)
[x^2 | x <- [1..5]]               -- quadrati
[x | x <- xs, x `mod` 2 == 0]     -- filtri / guardie
[(x,y) | x <- xs, y <- ys]        -- prodotto cartesiano
\end{haskellbox}

\subsubsection{Funzioni su Liste}
\begin{itemize}
    \item \textbf{Estrazione:} \code{head}, \code{tail}. \code{last}, \code{init}. \code{(!!)} per estrarre all'indice $n$
        \begin{haskellbox}{}
head [1, 2, 3]  -- risultato: 1
tail [1, 2, 3]  -- risultato: [2, 3]
last [1, 2, 3]  -- risultato: 3
init [1, 2, 3]  -- risultato: [1, 2]
[1, 2, 3] !! 1  -- risultato: 2
\end{haskellbox}{}

    \item \textbf{Ispezione:} \code{null}, \code{length}
\begin{haskellbox}{}
null []          -- risultato: True
length [1, 2, 3] -- risultato: 3
\end{haskellbox}{}

    \item \textbf{Taglio:} \code{take n}, \code{drop n}, \code{splitAt n} (tupla con take e drop).
\begin{haskellbox}{}
take 2 [1, 2, 3, 4]    -- risultato: [1, 2]
drop 2 [1, 2, 3, 4]    -- risultato: [3, 4]
splitAt 2 [1, 2, 3, 4] -- risultato: ([1, 2], [3, 4])
\end{haskellbox}{}

    \item \textbf{Generali:} \code{++} (concatenazione), \code{reverse}
\begin{haskellbox}{}
[1, 2] ++ [3, 4]  -- risultato: [1, 2, 3, 4]
reverse [1, 2, 3] -- risultato: [3, 2, 1]
\end{haskellbox}{}

    \item \textbf{Matematica/logica:} \code{sum}, \code{minimum}, \code{maximum}. \code{and}, \code{or}. \code{any p}, \code{all p} (predicati su almeno uno o tutti)
        \begin{haskellbox}{}
sum [1, 2, 3]       -- risultato: 6
minimum [5, 2, 9]   -- risultato: 2
maximum [5, 2, 9]   -- risultato: 9
and [True, False]   -- risultato: False
or [True, False]    -- risultato: True
any (> 2) [1, 2, 3] -- risultato: True (almeno uno è > 2)
all (> 0) [1, 2, 3] -- risultato: True (tutti sono > 0)
\end{haskellbox}{}

    \item \textbf{Liste annidate/multiple:} \code{concat} (appiattisce liste di liste), \code{zip} (unisce liste in coppie, fermandosi alla più corta)
\begin{haskellbox}{}
concat [[1, 2], [3, 4]]    -- risultato: [1, 2, 3, 4]
zip [1, 2] ['a', 'b', 'c'] -- risultato: [(1, 'a'), (2, 'b')]
\end{haskellbox}{}
\end{itemize}

\section{Funzion(al)i di Ordine Superiore}
I funzionali sono funzioni che prendono altre funzioni come argomento o le restituiscono come risultato. Favoriscono la \textbf{composizionalità}, permettendo di costruire programmi complessi tramite piccoli blocchi riutilizzabili.

\begin{itemize}
    \item \textbf{curryficazione:} si fonda sull'isomorfismo $A \times B \rightarrow C \cong A \rightarrow (B \rightarrow C)$. Le versioni curryficate permettono di passare un solo argomento alla volta, favorendo l'applicazione parziale.
    \item \textbf{$\eta$-regola:} è possibile omettere un parametro in una definizione se esso compare identico alla fine di entrambi i lati dell'equazione.
\end{itemize}

\begin{haskellbox}{}
-- composizione (.) : Equivale a f(g(x))
(.) f g x = f (g x)

-- curry e uncurry (Isomorfismo A x B -> C  < = = >  A -> (B -> C))
curry :: ((a, b) -> c) -> a -> b -> c
uncurry :: (a -> b -> c) -> (a, b) -> c

-- flip: inverte l'ordine dei parametri
flip :: (a -> b -> c) -> b -> a -> c
myFlip f a b = f b a

-- flip (:) [2, 3] 1
-- [1, 2, 3]

-- map e filter
map :: (a -> b) -> [a] -> [b]         -- proprieta': map (f . g) = map f . map g
filter :: (a -> Bool) -> [a] -> [a]

-- zipWith: applica f. binaria agli elementi di due liste
zipWith :: (a -> b -> c) -> [a] -> [b] -> [c]
\end{haskellbox}

\subsection{Famiglia dei \code{fold} (schemi di ricorsione)}
I \code{fold} astraggono il concetto di ricorsione strutturale sulle liste, sostituendo i costruttori della struttura dati (l'operatore \code{(:)} e la lista vuota \code{[]}) con funzioni e valori specifici. Se l'operatore \code{\#} è associativo e ha \code{e} come elemento neutro, allora $\code{foldl (\#) e} = \code{foldr (\#) e}$.

\begin{itemize}
    \item \textbf{\code{foldr} (right fold):} associa a destra; la riduzione avviene "al ritorno" dalle chiamate ricorsive (costruisce espressioni). 
\begin{itemize}
    \item sostituisce ricorsivamente il costruttore \code{(:)} con la funzione \code{f} data e \code{[]} con il valore base \code{z}.
        \item espansione: \code{foldr f z [x1, x2, x3]} diventa \code{x1 `f` (x2 `f` (x3 `f` z))}
        \item può terminare anche su liste infinite se la funzione \code{f} non è stretta nel suo secondo argomento (ovvero se non ha sempre bisogno di valutare la ricorsione per restituire un risultato).
\end{itemize}
        \begin{haskellbox}{}
-- somma con foldr
foldr (+) 0 [1, 2, 3]
-- valutazione: 1 + (2 + (3 + 0)) = 6

-- implementazione di map usando foldr
myMap f xs = foldr (\x acc -> f x : acc) [] xs

\end{haskellbox}

    \item \textbf{\code{foldl} (left fold):} associa a sinistra; simula un ciclo iterativo passando un accumulatore "in avanti"
\begin{itemize}
    \item valuta l'espressione partendo da sinistra e procedendo verso l'interno - valuta tutto solo alla fine della lista
        \item espansione: \code{foldl f z [x1, x2, x3]} diventa \code{((z `f` x1) `f` x2) `f` x3}
        \item \textit{nota bene:} In \code{foldl}, la funzione prende come primo argomento l'accumulatore, e come secondo l'elemento corrente della lista.
\end{itemize}
        \begin{haskellbox}{}
-- sottrazione con foldl
foldl (-) 0 [1, 2, 3]
-- valutazione: ((0 - 1) - 2) - 3 = -6

-- sottrazione con foldr per confronto
foldr (-) 0 [1, 2, 3]
-- valutazione: 1 - (2 - (3 - 0)) = 2

-- reverse usando foldl
myReverse xs = foldl (\acc x -> x : acc) [] xs
\end{haskellbox}

    \item \textbf{\code{foldr1} e \code{foldl1}:} varianti che non richiedono un valore iniziale esplicito, poiché usano il primo o l'ultimo elemento come base; \textit{danno errore su lista vuota}.
\begin{haskellbox}{}
-- utile quando non esiste un elemento neutro sensato (es. min/max assoluto)
myMinimum = foldr1 min
myMinimum [5, 3, 8] -- risultato: 3
-- myMinimum []     -- genera errore a runtime
\end{haskellbox}{}
\end{itemize}

\section{Alcuni esempi di stile funzionale}
Lo stile funzionale evita accumulatori ed indici, preferendo la manipolazione globale delle strutture dati.

\begin{itemize}
    \item \textbf{prodotto scalare:} può essere implementato in diversi modi (flessibilità dei funzionali):

\begin{haskellbox}{}
-- accoppia gli elementi, trasforma "*" in una funzione che accetta una tupla
-- applica la moltiplicazione agli elementi, e somma i prodotti
ps1 xs ys = sum (map (uncurry (*)) (zip xs ys))

-- crea una lista di funzioni parzialmente applicate
-- (da [5, 4] si ottiene [(5*), (4*)])
-- zApp/applyL prende la lista di funzioni ^ e la applica sulla seconda lista
ps2 xs ys = sum (applyL (map (*) xs) ys)

-- zipWith fa zip e map in un passo solo
ps3 xs ys = sum (zipWith (*) xs ys)
\end{haskellbox}

    \item controllare se una lista è ordinata:
\begin{haskellbox}{}
ordinata xs = and (zipWith (<=) xs (tail xs))
\end{haskellbox}

\item \textbf{ricerca su lista} (\code{findVPH}): si può implementare accoppiando la lista agli indici (\code{zip [0..] xs}) e filtrando per il valore cercato.
\end{itemize}

\chapter{Lambda Calcolo}
Ideato da Alonzo Church nel 1931, il $\lambda$-calcolo nasce per esplicitare il \textit{processo meccanico di calcolo} (sintassi) di una funzione, in contrasto con la visione matematica classica che vede la funzione solo come una relazione statica tra insiemi (semantica).

\section{Sintassi e Computazione}
La sintassi del $\lambda$-calcolo è minimalista ed è composta da tre regole:
\begin{itemize}
    \item \textbf{variabile:} $x, y, z \dots$
    \item \textbf{astrazione:} $\lambda x. M$ (corrisponde alla definizione di una funzione con parametro $x$ e corpo $M$)
    \item \textbf{applicazione:} $(M N)$ (corrisponde all'applicazione della funzione $M$ all'argomento $N$)
\end{itemize}

\begin{gframe}{Regole di Associazione}
    \begin{itemize}
        \item \textbf{l'applicazione associa a sinistra:} $F N_1 N_2 \dots N_n$ equivale a $(\dots (F N_1) \dots N_n)$
        \item \textbf{l'astrazione associa a destra:} $\lambda x_1 \dots x_n. M$ equivale a $\lambda x_1. (\lambda x_2. (\dots (\lambda x_n. M)\dots))$
    \end{itemize}
\end{gframe}

La computazione avviene tramite la \textbf{$\beta$-riduzione} (beta-regola): un termine della forma $(\lambda x.M) N$ si dice \textit{redex} e si riduce sostituendo tutte le occorrenze di $x$ in $M$ con il termine $N$.
Formalmente:
$$(\lambda x.M) N \rightarrow_\beta M[N/x]$$

Un termine che non contiene più $\beta$-redex è detto in \textbf{forma normale} e rappresenta un valore finale (il risultato della computazione).

\section{Variabili e Combinatori}
\begin{itemize}
    \item \textbf{Variabili libere ($FV$) e legate:} in $\lambda x. M$, la variabile $x$ è \textit{legata} (bound) all'astrazione. Qualsiasi variabile non legata è \textit{libera} (free).
    \item \textbf{$\alpha$-conversione (alfa-regola):} per evitare la cattura accidentale di variabili libere durante la sostituzione, si possono rinominare le variabili legate (es. $\lambda x.x \equiv \lambda y.y$). 
        \begin{itemize}
            \item è buona pratica seguire la ``convenzione di Barendregt'', che stabilisce che, all'interno di un termine lambda o di un insieme di termini, tutte le variabili legate devono avere nomi distinti tra loro e diversi da tutte le variabili libere presenti nel termine stesso
        \end{itemize} 
    \item  I termini tali che ($FV(M) = \emptyset$) ($\lambda$-termini chiusi), ovvero senza variabili libere sono chiamati \textbf{combinatori}
\end{itemize}

\subsection{Combinatori noti e strutture di controllo}
Il $\lambda$-calcolo puro non ha tipi di dato predefiniti: tutto (booleani, numeri, if-then-else) deve essere codificato tramite funzioni.

\begin{itemize}
    \item \textbf{Identità ($I$):} $I \equiv \lambda x.x$
    \item \textbf{Cancellatori (booleani):}
        \subitem \textbf{True ($K$):} $K \equiv \lambda x y. x$ (prende due argomenti e restituisce il primo)
        \subitem \textbf{False ($O$):} $O \equiv \lambda x y. y$ (prende due argomenti e restituisce il secondo)
    \item \textbf{If-Then-Else:} grazie alla definizione dei booleani, la struttura condizionale è banalmente codificata come un'applicazione: $if\ x\ then\ M\ else\ N \equiv x M N$. Se $x$ è True ($K$) selezionerà $M$, se è False ($O$) selezionerà $N$.
    \item \textbf{Compositori ($S$):} $S \equiv \lambda x y z. x z (y z)$ (si può dimostrare che $S K K \approx I$).
    \item \textbf{Duplicatori e divergenza:} 
        \subitem $\omega \equiv \lambda x.xx$ (auto-applicazione)
        \subitem $\Omega \equiv \omega \omega \rightarrow \omega \omega \rightarrow \dots$ rappresenta la \textbf{funzione indefinita}, il loop infinito (non tipabile)
\end{itemize}

\subsubsection{Numeri di Church e iteratori}
I \textbf{Numerali di Church} rappresentano i numeri naturali come funzioni di ordine superiore che \textit{iterano} un'operazione. Il numero $n$ è una funzione che prende una funzione (es. successore) $s$ e un valore di partenza (es. zero) $z$, e applica $s$ ad $z$ esattamente $n$ volte
$$\underline{n} \equiv \lambda s z. s(s(\dots(sz)\dots)) \quad [n\text{ applicazioni}]$$
\begin{itemize}
    \item $\underline{0} \equiv \lambda s z. z$ (equivale a \code{False})
    \item $\underline{3} \equiv \lambda s z. s(s(s(z)))$
\end{itemize}

Per calcolare il successore ($+1$) di un numero di Church, basta creare una funzione che applica $s$ una volta in più: $succ \equiv \lambda x s z. s(x s z)$ (o anche $\lambda x s z. xs (s z)$)

 
I numerali di Church non sono quindi solo dati, ma \textbf{iteratori intrinseci}.

\section{Ricorsione e Combinatori di Punto Fisso}
Per definire funzioni ricorsive complesse il $\lambda$-calcolo necessita di un principio di ricorsione generale. 

\begin{thmframe}{Combinatore di Punto Fisso}{}
 Nel $\lambda$-calcolo esiste un termine $Y$ (detto \textbf{Combinatore di Punto Fisso}) tale che, per ogni termine $M$, si ha $Y M \rightarrow M (Y M)$. 
 
 Il combinatore permette di trovare quindi il punto fisso di una funzione (quel valore $x$ tale che $f(x) = x$)
\end{thmframe}

Due combinatori di punto fisso sono:
\begin{itemize}
    \item \textbf{Combinatore di Turing ($Y_T$):} definito come $Y_T = \theta \theta$ dove $\theta \equiv \lambda x y. y(x x y)$.
    \item \textbf{Combinatore di Curry ($Y_C$):} definito come $Y_C \equiv \lambda f. (\lambda x.f(xx))(\lambda x.f(xx))$. Se applicato all'identità $I$, questo combinatore riduce inesorabilmente a $\Omega$ (divergenza).
\end{itemize}

\begin{thmframe}{Tesi di Church-Turing}{}
 Ogni funzione calcolabile è $\lambda$-definibile.
\end{thmframe}



\end{document}
